% =============================================================================
% Section 2 — Mathematical Formulation
% PINN-ONB01 Phase 1 manuscript (IJHMT submission)
% Target length: ~800 words (equations and tables excluded)
% =============================================================================

\section{Mathematical Formulation}\label{sec:formulation}

% -----------------------------------------------------------------------------
\subsection{Governing equations and boundary conditions}\label{subsec:governing}%
\label{subsec:setup}\label{subsec:pde}\label{subsec:bc}
% -----------------------------------------------------------------------------

We consider a horizontal heater of thickness $L$ immersed in a quiescent
liquid pool at saturation pressure. The $z$-axis is taken normal to the
heating surface, with origin at the bottom face. A uniform heat flux
$q''_{\mathrm{applied}}$ is imposed at $z=0$, while the fluid-facing face
is at $z=L$. Because the lateral extent exceeds $L$ by at
least two orders of magnitude, the temperature field is treated as
one-dimensional. The 1D approximation is further justified by the
Biot number $\mathrm{Bi} = h_{nc}\,L/k_s$. For the dataset considered
here ($h_{nc} \lesssim \SI{2000}{W/(m^2 K)}$, $L \le \SI{5}{mm}$,
copper or stainless heaters with $k_s = 16$--$\SI{400}{W/(m\,K)}$),
$\mathrm{Bi}$ stays in the $10^{-4}$--$10^{-2}$ range, so radial
gradients within the heater remain negligible relative to those in
the thickness direction. The following assumptions are adopted:
(i) steady-state conduction, since ONB lies on the steady single-phase
branch; (ii) negligible volumetric heat generation; (iii) adiabatic lateral
boundaries; (iv) uniform, temperature-independent solid conductivity $k_s$;
and (v) laminar natural convection above $z=L$, closed by a Nusselt-type
correlation evaluated at $T(L)$.
The same framework applies to large-diameter horizontal cylinders; for
brevity, the derivation is presented for the flat-plate geometry.

Under these assumptions, the temperature distribution $T(z)$ within the
heater is governed by the steady, source-free, one-dimensional heat
equation
%
\begin{equation}\label{eq:laplace}
\frac{\mathrm{d}^{2} T}{\mathrm{d} z^{2}} \;=\; 0
\qquad \text{for } z \in [0,L].
\end{equation}
%
With $q'''=0$, \cref{eq:laplace} admits the exact linear solution
$T(z) = T(0) - (q''_{\mathrm{applied}}/k_s)\, z$ once
\cref{eq:neumann_wall} is enforced. This closed form does not supplant
the PINN; rather, it serves three roles: (i) an analytical reference for
Level-1 verification, (ii) a regularizing signal during warm-up, and (iii)
the basis for the thermal-boundary-layer thickness $\delta_t$ that enters
\cref{eq:hsu_rc_range}.

Equation~\eqref{eq:laplace} is closed by a Neumann condition at $z=0$ and
a Robin condition at $z=L$, given respectively by the prescribed applied
flux,
%
\begin{equation}\label{eq:neumann_wall}
-k_s\,\left.\frac{\mathrm{d}T}{\mathrm{d}z}\right|_{z=0}
\;=\; q''_{\mathrm{applied}},
\end{equation}
%
while conjugate coupling with the pool is enforced via Newton's law of
cooling, using the natural-convection coefficient $h_{nc}$:
%
\begin{equation}\label{eq:nc_bc}
-k_l\,\left.\frac{\mathrm{d}T}{\mathrm{d}z}\right|_{z=L}
\;=\; h_{nc}\,\bigl(T_L - T_{\infty}\bigr),
\end{equation}
%
where $k_l$ denotes the liquid thermal conductivity,
$T_L \equiv T(z{=}L)$ the surface-side temperature, $T_{\infty}$ the
bulk-pool temperature, and
$\Delta T_{\mathrm{sub}} \equiv T_{\mathrm{sat}} - T_{\infty}$ the
subcooling. For an
upward-facing heated surface, $h_{nc}$ follows
McAdams~\citep{incropera,lienhard1981}:
%
\begin{equation}\label{eq:nu_nc}
\mathrm{Nu}_L \;=\;
\begin{cases}
0.54\,\mathrm{Ra}_L^{1/4}, & 10^{4} \le \mathrm{Ra}_L \le 10^{7} \quad\text{(laminar)},\\[2pt]
0.15\,\mathrm{Ra}_L^{1/3}, & 10^{7} \le \mathrm{Ra}_L \le 10^{11} \quad\text{(turbulent)},
\end{cases}
\end{equation}
%
with $\mathrm{Ra}_L = g\,\beta_l\,(T_L - T_{\infty})\,L^{3}/(\nu_l\,\alpha_l)$.
All cases in the present dataset fall within the laminar branch; the
turbulent branch is retained for completeness. When only the bulk-pool
temperature is reported, the Dirichlet substitute $T(L) = T_{\infty}$ is
applied in the data-loss term, whereas the PDE residual continues to
enforce \cref{eq:nc_bc}.

% -----------------------------------------------------------------------------
\subsection{Hsu nucleation criterion}\label{subsec:hsu}
% -----------------------------------------------------------------------------

Incipience is modeled following Hsu's criterion~\citep{hsu1962}, which
identifies the range of cavity radii that simultaneously satisfy the
mechanical-equilibrium and thermal-boundary-layer conditions required for
vapor-embryo growth. With a linear
profile of thickness $\delta_t = k_l\,\Delta T_{\mathrm{wall}}/q''$, the
active band is bounded by
%
\begin{equation}\label{eq:hsu_rc_range}
r_{c,\,\min/\max} \;=\;
\frac{\delta_t}{2\,C_1}\,
\left[\,1 \mp \sqrt{\,1 - \frac{8\,C_2\,\sigma\,T_{\mathrm{sat}}}{\delta_t\,\Delta T_{\mathrm{wall}}\,\rho_v\,h_{fg}}}\,\right],
\end{equation}
%
with $C_1 = (1+\cos\theta)/\sin\theta$ and
$C_2 = (1+\cos\theta)/\sin^{2}\theta$; $\sigma$ is the surface tension,
$\rho_v$ the saturated vapor density, and $h_{fg}$ the latent heat. For
$\theta = 90^{\circ}$, the default in the correlation-comparison module,
$C_1 = C_2 = 1$. Requiring the discriminant to be non-negative yields the
necessary nucleation condition
%
\begin{equation}\label{eq:hsu_discriminant}
\delta_t\,\Delta T_{\mathrm{wall}}
\;\ge\; \frac{8\,C_2\,\sigma\,T_{\mathrm{sat}}}{\rho_v\,h_{fg}}.
\end{equation}
%
Equality collapses the band to a single radius. The corresponding minimum
superheat is
%
\begin{equation}\label{eq:hsu_dt_onb}
\Delta T_{\mathrm{ONB,\,Hsu}}
\;=\; \sqrt{\,\frac{8\,C_2\,\sigma\,T_{\mathrm{sat}}\,q''}{k_l\,\rho_v\,h_{fg}}\,}.
\end{equation}
%
Equation~\eqref{eq:hsu_dt_onb} provides the analytical baseline used in
Section~\ref{sec:results} and is embedded as a soft penalty within the
composite loss. The $\theta$-dependent prefactor allows the network to
inherit a wettability sensitivity that scalar correlations cannot
capture.

% -----------------------------------------------------------------------------
\subsection{Non-dimensionalization}\label{subsec:nondim}\label{subsec:dimless}
% -----------------------------------------------------------------------------

The system is recast in three reference quantities to remove
fluid-specific scales and improve loss conditioning. The capillary
length is $L_c$. The latent-heat temperature scale is
$\Delta T_{\mathrm{ref}}$. A Zuber-type heat-flux scale is
$q_{\mathrm{ref}}$, with
%
\begin{equation}\label{eq:Lc}
L_c \;=\; \sqrt{\,\frac{\sigma}{g\,(\rho_l - \rho_v)}\,},
\end{equation}
%
\begin{equation}\label{eq:dT_ref}
\Delta T_{\mathrm{ref}} \;=\; \frac{h_{fg}}{c_{p,l}},
\end{equation}
%
\begin{equation}\label{eq:q_ref}
q_{\mathrm{ref}} \;=\; h_{fg}\,\rho_v^{1/2}\,\bigl[\sigma\,g\,(\rho_l - \rho_v)\bigr]^{1/4},
\end{equation}
%
with $h_{\mathrm{ref}} = q_{\mathrm{ref}}/\Delta T_{\mathrm{ref}}$.
\cref{eq:q_ref} is the Zuber CHF scale with $K=1$~\citep{lienhard1981}.
The engineering CHF is recovered upon multiplication by $K \approx 0.131$.
Dimensionless variables are defined as
%
\begin{equation*}
z^{\ast} = z/L_c, \qquad
T^{\ast} = (T-T_{\mathrm{sat}})/\Delta T_{\mathrm{ref}}, \qquad
q''^{\ast} = q''/q_{\mathrm{ref}}, \qquad
r_c^{\ast} = r_c/L_c.
\end{equation*}
%
For water at atmospheric pressure, $L_c \approx \SI{2.5}{mm}$,
$\Delta T_{\mathrm{ref}} \approx \SI{2245}{K}$, and
$q_{\mathrm{ref}} \approx \SI{8.46}{MW/m^2}$. These yield a CHF of
$\approx \SI{1.1}{MW/m^2}$, consistent
with~\citet{zuber1959,lienhard1981}. These reference quantities are
recomputed for each fluid using the routine of Section~\ref{sec:data}, so
that cross-fluid generalization is performed in a common non-dimensional
frame.

Four groups summarize the physics and serve as auxiliary PINN inputs. The
wall-superheat Jacob number
%
\begin{equation}\label{eq:ja_sup}
\mathrm{Ja} \;=\;
\frac{\rho_l\,c_{p,l}\,\Delta T_{\mathrm{wall}}}{\rho_v\,h_{fg}}
\end{equation}
%
controls bubble growth. A value of $\mathrm{Ja}=\mathcal{O}(1)$ marks
inertia-controlled growth. The subcooled counterpart
$\mathrm{Ja}_{\mathrm{sub}} = \rho_l c_{p,l}\,\Delta T_{\mathrm{sub}}/(\rho_v\,h_{fg})$
weights cap condensation. The Bond number
%
\begin{equation}\label{eq:bond}
\mathrm{Bo} \;=\;
\frac{(\rho_l - \rho_v)\,g\,L^{2}}{\sigma}
\end{equation}
%
equals unity at $L=L_c$, motivating \cref{eq:Lc}. The liquid Prandtl
number $\mathrm{Pr}_l = \mu_l\,c_{p,l}/k_l$ sets the relative thickness
of viscous and thermal layers, entering $h_{nc}$ through
$\mathrm{Ra}_L = \mathrm{Gr}_L\,\mathrm{Pr}_l$. Combined with the surface
descriptors of Section~\ref{sec:data}, these four numbers form the
complete dimensionless input to the PINN of
Section~\ref{sec:architecture}.
